\section{Prototype}
\label{sec:prototype}

\begin{figure}[t]
    \centering
    \begingroup
    \setlength{\fboxsep}{3pt}
    \setlength{\tabcolsep}{2pt}
    \resizebox{\linewidth}{!}{%
    \begin{tabular}{c c c c c c c}
        \fbox{\parbox[c][0.52cm][c]{1.45cm}{\centering\scriptsize XIAO\\microphone}}
        & $\rightarrow$ &
        \fbox{\parbox[c][0.52cm][c]{1.45cm}{\centering\scriptsize Audio\\window}}
        & $\rightarrow$ &
        \fbox{\parbox[c][0.52cm][c]{1.55cm}{\centering\scriptsize Log-mel\\features}}
        & $\rightarrow$ &
        \fbox{\parbox[c][0.52cm][c]{1.55cm}{\centering\scriptsize Integer\\inference}} \\
        \multicolumn{7}{c}{$\downarrow$} \\
        \fbox{\parbox[c][0.52cm][c]{1.45cm}{\centering\scriptsize Intent +\\confidence}}
        & $\rightarrow$ &
        \fbox{\parbox[c][0.52cm][c]{1.45cm}{\centering\scriptsize Timing\\fields}}
        & $\rightarrow$ &
        \fbox{\parbox[c][0.52cm][c]{1.55cm}{\centering\scriptsize Safety-state\\bridge}}
        & $\rightarrow$ &
        \fbox{\parbox[c][0.52cm][c]{1.55cm}{\centering\scriptsize Gated Tello\\SDK command}}
    \end{tabular}%
    }
    \endgroup
    \caption{Prototype data path. The board performs local audio capture, feature extraction, and integer inference, then exposes intent-event fields to a safety-state bridge that gates Tello SDK command dispatch.}
    \label{fig:prototype_pipeline}
\end{figure}


\subsection{Embedded Voice-Event Source}

The prototype instantiates the voice layer as a local event source near the UAV. A XIAO ESP32-S3 Sense captures onboard audio, computes the acoustic frontend, invokes the compact recognizer, and emits an intent event with confidence and timing fields. The board is not a flight controller in this design. Its responsibility is to turn rotor-noisy speech into a structured event record that can be checked by a separate control boundary.

Table~\ref{tab:prototype_interfaces} summarizes the prototype interfaces and the evidence each interface is meant to support. The table is a placeholder for the implementation description; measured latency, throughput, and success rates belong in Section~\ref{sec:evaluation}.

\begin{table}[t]
    \centering
    \caption{Prototype scope. The prototype demonstrates local voice-event generation while keeping control decisions outside the recognizer.}
    \label{tab:prototype_scope}
    \begin{tabular}{p{1\linewidth}}
        \toprule
        Prototype demonstrates  \\
        \midrule
        The UAV-side device can listen near the drone and produce intent events locally.\\
        The recognizer output can be packaged as an event with confidence and timing.  \\
        BLE/USB reporting can carry events to a host-side bridge.  \\
        The bridge can log event, state, decision, and acknowledgement fields. \\
        \bottomrule
    \end{tabular}
\end{table}

\subsection{Runtime Event Interface}

The embedded event should contain the minimum fields needed for downstream inspection: event label, confidence, timestamp or sequence number, timing breakdown, and raw-output diagnostics. The bridge can then associate each event with the current UAV state, selected boundary decision, acknowledgement, timeout, or rejection. This interface keeps recognition and control separate: the board reports what it heard and how confident it was, while the bridge decides whether the event is ignored, held, escalated, or forwarded.

The current runtime candidate is used to test whether this interface can execute on the ESP32 path under the model's operator and memory constraints. Those measurements should be reported as deployment feasibility and response-time evidence. They should not be described as labeled real-world semantic accuracy unless the same events are paired with labeled acoustic ground truth.
% TODO(results): Add runtime candidate provenance and timing table in Section~\ref{sec:evaluation}, not here.
% TODO(boundary): Keep first-window latency, steady-state output period, and pure model invocation as separate measurements.

\subsection{Control-Bridge Boundary}

The prototype keeps the control bridge outside the board for the current draft. This host-mediated boundary is deliberate: it lets the board remain a voice-event source while the host records policy checks, command eligibility, acknowledgement behavior, and manual override status. Movement events therefore remain requests until a downstream policy authorizes action. Emergency events can be prioritized or escalated, but they still pass through the boundary. Unknown and low-confidence events should be logged and rejected or mapped to no action.

This boundary is also where the remaining system validation must occur. A dry run can test message formats and state transitions, and a grounded or no-propeller setup can test bridge behavior without claiming flight validation. Live-flight safety claims require stronger evidence than the current prototype placeholder.
% TODO(validation): Add no-propeller or grounded bridge logs before describing bridge behavior as validated.
% TODO(validation): Add gate/buffer false-trigger results before using the trigger gate as part of the safety argument.